\chapter{武汉大学2016年博士学位论文撰写及印制规格的规定}


{ {\heiti 说明}: \kaishu
 这些规定偶有变动.
 请仔细查阅当年拿到的《武汉大学申请博士学位有关规定》那个红色小本子.
 若发现不一致的地方, 请您发邮件给我: huangzh@whu.edu.cn.}

\bigskip


[原文地址: \url{http://www.gs.whu.edu.cn/index.php/index-view-aid-7956.html}.]
\bigskip

国务院学位办文件规定：研究生学位论文撰写和编印必须符合国家标准局统一编制格式，
并将博士学位论文提交北京图书馆和北京科技信息研究所收入国家级图书编目，
以便及时向社会提供查阅，促进国内外学术交流. 根据文件精神结合我校学位论文归档的要求，
对我校博士学位论文印制规格作出如下规定：

{\heiti 一、论文装订格式的排列顺序}

（一）封面（格式和要求见附页一）

（二）论文英文题目（专用一页纸，上方为题目用宋体2号字，下方为研究生姓名宋体4号字、外文专业应有中文题目）

（三）郑重声明（格式和要求见附页二）

（四）学位论文使用授权书（格式和要求见附页三）

（五）博士生自认为的论文创新点

（六）目录（格式和要求见附页四）

（七）中文摘要

（八）英文摘要

（九）引言

（十）正文（格式和要求见附页五）

（十一）中外文参考文献（格式和要求见附页六）

（十二）攻博期间发表的与学位论文相关的科研成果目录（格式参见附页六：参考文献排列格式）

（十三）后记/致谢

{\heiti 二、论文印制规格及要求}

（一）论文用A4张（210×297mm）标准大小的白纸打印；\colorbox{yellow}{正文部分双面印制}.

（二）论文在打印时，纸张四周留足空白边缘，每页上方（天头）和左侧（订口）应分别留边25mm以上，
下方（地脚）和右侧（切口）分别留边25mm以上.

{\heiti 三、论文封面格式}

（一）分类号. 必须在封面左上角注明分类号，一般应注明《中国图书资料分类法》的类号，
同时尽可能注明《国际十进分类法UDC》的类号.

（二）密级. 论文必须按国家规定的保密条例在右上角注明密级（如系公开型论文则可不注明密级）.

（三）编号. 武汉大学编号为10486，标注在封面右上角.

（四）论文题目. 题目必须用楷体标准一号字标注于明显的位置，应是集中概括论文最重要的内容，
一般不超过20个字，以有助于选定关键词和编制题录. 题目不能用缩略词，首字母缩写字、字符、
代号和公式等，题目语意未尽，可用副标题补充说明. 外语专业的论文题目一般采用英文，英文题目不宜超过10个实词.

（五）论文作者姓名.

（六）论文指导教师姓名. 指导教师姓名必须是填写当年被学校批准招收博士生的教师.

（七）专业名称. 专业名称必须是我校已有学位授予权的学科专业，并按国家颁布的学科专业目录中一级学科或二级学科名称印制.

（八）书脊（专指博士学位论文）. 书脊上应用仿宋体四号字于上方标明论文题目，下方注明研究生姓名.

（九）论文封面统一用120克铜版纸，封面底色为白色.

（十）论文封面格式要求详见附三页.

{\heiti 四、论文英文题目}

论文英文题目专用一页纸，``英文题目''用宋体二号字，其下``研究生姓名''用宋体四号字；外语专业应为中文题目.

{\heiti 五、郑重声明}


``郑重声明''用黑体小二号字，内容用宋体四号字.

为了加强学风、学术道德建设，规范学术行为，提高学位论文质量，确保学位授予的权威性、严肃性，学校对学位论文撰写特别强调以下几点：

（一）凡申请学位人员须对自己的学位论文负责，在提交的学位论文的英文题目后页（中文摘要前页）增设一页书面声明，即``郑重声明''.

（二）学位论文中的引证、引述处须注明出处.

（三）学位论文后所附参考文献，必须是申请学位人员真正阅读和参考过的文献.

（四）合作科研及成果，应在致谢或后记中有相关说明，避免产权纠纷.

（五）学位论文中没有郑重声明的，不能参加论文答辩.

{\heiti 六、中文摘要}

``中文摘要''用黑体小二号字，内容用宋体小四号字，页码用罗马数字单独编排，并标注在每页页脚中部.

{\heiti 七、英文摘要}

``英文摘要''用加粗Times New Roman小二号字，内容用Times New Roman小四号字，页码续接中文摘要的页码.

{\heiti 八、论文关键词}

每篇论文必须选取3--5个以上中、英文关键词，排在其论文摘要的左下方，用黑体小四号字.

{\heiti 九、目录}

目录是论文的提纲，也是论文组成部分的小标题.
排列顺序是：1、中文摘要 2、英文摘要 3、引言 4、正文章节 5、中外文参考文献 6、攻博期间发表的科研成果目录
7、后记（可不要此项）. 并对每项标明页码.

{\heiti 十、引言（绪论）}

论文的页码由引言（绪论）的首页开始，作为第1页，并为右页，一律用阿拉伯数字连续编排页码，必须统一标注在每页页脚中部.

{\heiti 十一、正文}

正文是学位论文的核心部分，必须由另页开始，一级标题之间换页，二级标题之间空行；内容一律用宋体小四号字，
字间距设置为标准字间距，行间距设置为最小值20磅，各章、节应有序号.

{\heiti 十二、参考文献}

参考文献用黑体四号字，内容用宋体五号字.

{\heiti 十三、学位论文印刷份数}

由培养单位根据需要决定印刷份数. 在学位论文定稿后，可先印刷数份给论文评阅人和答辩委员，
然后根据论文评阅人和答辩委员对论文的意见进行修改后，才能正式印刷，提交学校存档.

{\heiti 十四、论文制作时须加页眉}


页眉从中文摘要开始至论文末, 偶数页码内容为``武汉大学博士学位论文'', 奇数页码内容为学位论文题目.

